\chapter{Configuration Key Reference}\label{app:config}

This appendix lists the configuration keys consumed by the engine, grouped by their YAML
section. It is a reference companion to the chapters: each key is described in context in
the relevant chapter, and here we collect them with their types and defaults. The keys and
defaults follow the engine's configuration module.

\paragraph{Load rules.} A configuration must contain all five top-level sections
(\conf{mode}, \conf{numerics}, \conf{mesh}, \conf{data}, \conf{time}); a section that is
present but empty must be written as \code{\{\}}. Keys that are misspelled or of the wrong
case are silently ignored and the default is used --- there is no typographic protection ---
so a key that appears to have no effect should be checked against this table. Relative paths of the mesh,
data, seed and output files are resolved against the directory of the configuration file;
\conf{numerics.shared\_fields\_dir} is taken as written. Many values are validated when the file is
loaded --- the integrator, the interpolation methods \conf{uv\_interp}, \conf{w\_interp} and
\conf{sample\_interp}, the diffusion, beaching (when enabled), periodic, adaptive-scheme, grid-hash, rotation
and HEALPix keys, \conf{w\_positive}, \conf{uv\_vertical}, \conf{input\_rotation} and the stamp keys of
\conf{data} --- and an invalid value there raises an error immediately; others, such as
\conf{max\_levels}, \conf{time\_averaging} and the \conf{time} values, are checked when the run starts or
when first used. Unknown keys inside the diffusion and
beaching blocks are ignored like any other.

\setlength{\tabcolsep}{4pt}

\section*{\texttt{mode} --- tracking mode and variable names}
\begin{longtable}{@{}>{\ttfamily\footnotesize}p{0.295\linewidth} >{\footnotesize}p{0.075\linewidth} >{\footnotesize}p{0.095\linewidth} >{\footnotesize}p{0.37\linewidth}@{}}
\toprule
\normalfont Key & Type & Default & Description \\
\midrule \endhead
Fdrift2D & bool & \code{true} & Kept for older configurations; not read by the engine. \\
Fdrift3D & bool & \code{false} & Three-dimensional tracking with vertical velocity; alone decides between 2-D and 3-D. \\
time\_dir & float & \code{1.0} & Integration direction: \code{1.0} forward, \code{-1.0} backward; any other value is refused. \\
Vpuls & bool & \code{false} & FESOM-C 2-D only: add the pulsation velocity to the mean flow. \\
var\_u, var\_v & str & auto & Names of the horizontal velocity variables; set to \code{unod}/\code{vnod} for node-layout FESOM2 velocities (requires \conf{uv\_interp: barycentric}). \\
var\_w & str & auto & Name of the vertical velocity (3-D). FESOM2: if unresolved, the run proceeds in 2-D silently; FESOM-C refuses \code{w\_cart}. \\
w\_positive & str & \code{auto} & Sign convention of the vertical velocity in the file: \code{auto}, \code{up} or \code{down} (\cref{sec:vert-wsign}). \\
var\_eta & str & auto & Name of the sea-surface height. \\
uv\_vertical & str & \code{layer} & Horizontal velocity in depth: \code{layer} (constant within each layer) or \code{linear} (linear between the layer centres; a stencil member whose bed lies above the particle is a wall; \cref{sec:vert-uvdepth}). \code{linear} needs 3-D tracking (refused at load otherwise), a z-level or sigma mesh, \code{euler}, \code{rk2} or \code{rk4}, and \conf{uv\_interp} \code{barycentric} for node layout or \code{leastsq}, \code{edge\_leastsq} or \code{barycentric} for element layout; FESOM2 element data need the element bottoms and no ice-shelf cavities (\cref{sec:mesh-elem-bottoms}). Refused for \code{fesom14}, whose velocities on the levels are always linear in depth. Recorded as the output attribute \code{uv\_vertical = "linear"}; a resume must keep it (\cref{sec:out-restart}). \\
sigma\_layer & int & surface & 2-D layered input: which vertical layer to track (\code{fesom14}: a 0-based kept level). \\
sigma\_mean & bool & \code{false} & FESOM-C and \code{fesom14} 2-D: thickness-weighted vertical mean instead of one layer. \\
max\_levels & int & all & Keep the first $N$ vertical interfaces from the surface; must be $\ge 2$ in 3-D. \\
sample\_vars & list & \code{[]} & Node-based scalar fields to sample along trajectories (e.g.\ \code{temp}, \code{salt}). \\
track\_failure\_reason & bool & \code{false} & Record why each particle died (\cref{sec:life-reason}). \\
\bottomrule
\end{longtable}

\section*{\texttt{numerics} --- step, integrator, interpolation}
\begin{longtable}{@{}>{\ttfamily\footnotesize}p{0.295\linewidth} >{\footnotesize}p{0.075\linewidth} >{\footnotesize}p{0.095\linewidth} >{\footnotesize}p{0.37\linewidth}@{}}
\toprule
\normalfont Key & Type & Default & Description \\
\midrule \endhead
dt & float (s) & auto & Integration step. Defaults to the snapshot interval. In the discrete mode it must equal the record spacing, and irregular record times (calendar months) are refused --- use \conf{temporal\_interp}, where it may be smaller. \\
file\_dt & float (s) & auto & Time between stored snapshots; auto-detected from the time coordinate. When the axis cannot be decoded it defines $t_i=i\,\Delta t$ (\cref{sec:input-axis}). \\
integrator & str & \code{euler} & \code{euler} / \code{rk2} / \code{rk4} / \code{rk45} (\cref{ch:integration}). \\
rk45\_tol\_m & float (m) & \code{1.0} & Adaptive scheme: error tolerance per substep, $>0$ (\cref{sec:int-rk45}). \\
rk45\_h\_min & float (s) & $10^{-4}\dt$ & Adaptive scheme: smallest substep, $>0$. \\
output\_freq & int & \code{1} & Record positions every $N$ steps, counted from the start of the run. \\
temporal\_interp & bool & \code{false} & Reconstruct velocity in time between snapshots; permits sub-stepping (\cref{ch:temporal}). \\
temporal\_reconstruction & str & \code{linear} & \code{linear} or \code{pchip} (monotone cubic); active only with \conf{temporal\_interp}. \\
velocity\_float32 & bool & \code{false} & Hold velocity in single precision; halves snapshot memory. \\
prefetch\_size & int & \code{1} & Snapshots read per bulk read-ahead; forced to 1 under cyclic forcing and with the shared cache. \\
shared\_fields\_dir & str & \code{null} & Directory of the node-shared snapshot cache (\cref{sec:input-shared}). \\
shared\_fields\_generation & str & \code{null} & Namespace generation of the cache; else \code{KELP\_SHARED\_FIELDS\_GENERATION}, else the batch job. \\
checkpoint\_freq & int & \code{0} & Write a restart checkpoint every $N$ steps (\code{0} = off); needs an output path. \\
output\_flush\_every & int & \code{0} & NetCDF only: flush every $N$ records so the file is readable mid-run. \\
uv\_interp & str & \code{leastsq} & \code{leastsq} (element centres), \code{edge\_leastsq} (containing element and edge neighbours) or \code{barycentric} (vertex blend; required for node velocities) (\cref{ch:hvel}). \\
edge\_leastsq\_bc & str & \code{freeslip} & Wall ghost of \code{edge\_leastsq}: \code{freeslip} or \code{noslip}. \\
elem\_search & str & \code{local} & Containing element: \code{local} (search each step) or \code{walk} (track across edges; \cref{sec:loc-walk}). \\
w\_interp & str & \code{leastsq} & 3-D vertical-velocity horizontal combine: \code{nearest}/\code{leastsq}/\code{barycentric}/\code{idw}. \\
sample\_interp & str & \code{leastsq} & Horizontal method for sampled scalars (same options). \\
Dmin & float (m) & \code{0.5} & Minimum water-column thickness (sigma meshes only; \cref{sec:vert-clamp}). \\
z0b\_min & float (m) & \code{0.1} & Bottom clamp offset (sigma meshes only). \\
Dch & float & \code{1.0} & Reserved; currently unused. \\
periodic\_lon, periodic\_lat & [lo, hi] & \code{null} & Periodic windows for channel domains (\cref{sec:life-periodic}). \\
diffusion & map & off & Stochastic-diffusion sub-configuration (below). \\
beaching & map & off & Beaching sub-configuration (below). \\
\bottomrule
\end{longtable}

\section*{\texttt{numerics.diffusion}}
\begin{longtable}{@{}>{\ttfamily\footnotesize}p{0.295\linewidth} >{\footnotesize}p{0.075\linewidth} >{\footnotesize}p{0.095\linewidth} >{\footnotesize}p{0.37\linewidth}@{}}
\toprule
\normalfont Key & Type & Default & Description \\
\midrule \endhead
enabled & bool & \code{false} & Master switch; when off, byte-identical to no diffusion. \\
scheme & str & \code{uniform} & \code{uniform} / \code{em} (Euler--Maruyama) / \code{m1} (Milstein). Only \code{uniform} for the horizontal kick in 3-D. \\
kh\_zonal, kh\_meridional & float & \code{0.0} & Constant diffusivities (\code{uniform} only), m$^2$/s. \\
var\_kh\_zonal, var\_kh\_meridional & str & auto & Dataset diffusivity variable names (\code{em}/\code{m1}). \\
dres\_m & float (m) & auto & Central-difference offset for the diffusivity gradient. \\
seed & int & \code{0} & Counter-based RNG seed (\cref{sec:diff-rng}); restart-safe. \\
composition & str & \code{split} & How the kick composes with the step: \code{split}/\code{combined}/\code{split\_rk}. \\
kh\_eval\_position & str & \code{end} & \code{split\_rk} only: sample diffusivity at \code{start} or \code{end} of step. \\
az\_enabled & bool & \code{false} & Vertical random walk (\cref{sec:diff-vertical}); needs \conf{enabled} and the \code{uniform} scheme. \\
az\_const & float & \code{0.0} & Constant vertical diffusivity, m$^2$/s, $\ge 0$. \\
var\_az & str & \code{null} & Vertical diffusivity variable on the interfaces (FESOM2 z-level). \\
\bottomrule
\end{longtable}

\section*{\texttt{numerics.beaching}}
\begin{longtable}{@{}>{\ttfamily\footnotesize}p{0.295\linewidth} >{\footnotesize}p{0.075\linewidth} >{\footnotesize}p{0.095\linewidth} >{\footnotesize}p{0.37\linewidth}@{}}
\toprule
\normalfont Key & Type & Default & Description \\
\midrule \endhead
enabled & bool & \code{false} & Master switch (\cref{sec:life-beach}). \\
criteria & list & \code{[depth]} & Any of \code{depth}, \code{coastal\_element}, \code{on\_failure}, \code{stochastic}; combined inclusively. \\
depth\_min & float (m) & \code{0.3} & Strand when the column height falls below this. \\
use\_eta & bool & \code{false} & Include sea level in the column height (needs an ssh field). \\
dry\_mask\_var & str & none & Node variable marking dry cells ($>0.5$); strands (with the \code{depth} criterion) and blocks resuspension. \\
zone\_depth & float (m) & auto & Near-shore zone for stochastic stranding; defaults to twice \conf{depth\_min}. \\
t\_beach & float (s) & \code{86400} & Stochastic stranding timescale: $P=1-\exp(-|\dt|/t_{\text{beach}})$ per step. \\
resuspension & str & \code{never} & One of \code{never}, \code{auto}, \code{stochastic}, or \code{auto\_or\_stochastic}. \\
rewet\_hysteresis & float (m) & \code{0.1} & Margin above \conf{depth\_min} to count as re-wetted. \\
min\_wet\_steps & int & \code{1} & Consecutive wet steps before automatic release. \\
t\_resuspend & float (s) & none & Stochastic-release timescale (required for stochastic resuspension). \\
seed & int & none & Beaching RNG seed; defaults to the diffusion seed. \\
\bottomrule
\end{longtable}

\section*{\texttt{mesh} --- mesh source}
\begin{longtable}{@{}>{\ttfamily\footnotesize}p{0.295\linewidth} >{\footnotesize}p{0.075\linewidth} >{\footnotesize}p{0.095\linewidth} >{\footnotesize}p{0.37\linewidth}@{}}
\toprule
\normalfont Key & Type & Default & Description \\
\midrule \endhead
type & str & \code{fesomc} & \code{fesom2} / \code{fesom14} / \code{fesomc} / \code{nemo} / \code{healpix}; selects the reader, conventions, and vertical coordinate. \code{fesom14} needs \conf{uv\_interp: barycentric} (\cref{sec:mesh-fesom14}). \\
meshpath & str & required & Path to the mesh file or directory (\cref{ch:mesh}). \code{fesom14}: an ASCII mesh directory (it holds \code{nod3d.out}), or level-wise data with cell polygons --- a file, or a directory whose first \code{*.nc} file holding \conf{mode.var\_u} is used. \\
input\_rotation & [3 floats] & \code{null} & \code{fesom14} only (refused at load for every other type): the Euler angles in degrees, from the run's namelist, of the rotated frame an ASCII mesh is stored in; rotated back to geographic on load. Refused for a mesh from cell polygons; independent of \conf{rotation}. \\
r\_earth & float (m) & family & Earth radius for mesh-derived areas; defaults per model. \\
rotation & varies & \code{null} & FESOM2, FESOM~1.4 and HEALPix: \code{null} (geographic), \code{fesom2}, or an explicit list of Euler angles (\cref{ch:frame}). \\
exact\_nearest & bool & \code{true} & Prove every nearest-node anchor (\cref{sec:loc-exact}). \\
grid\_hash\_cell\_rule & str & \code{occu-} \code{pancy} & Choose the hash cell by measurement (\code{occupancy}), or keep the reader's (\code{legacy}) (\cref{sec:loc-cell}). \\
grid\_hash\_cell\_size & float (deg) & \code{null} & Explicit hash cell; overrides the rule. \\
healpix\_zoom & int & \code{null} & HEALPix: zoom level when the store carries no grid attributes. \\
healpix\_geometry & str & \code{dual} & HEALPix: \code{dual} (nodes at pixel centres) or \code{pixel} (elements are pixels; 2-D). \\
crop\_bbox & [W, E, S, N] & \code{null} & HEALPix: crop to a box; $W>E$ crosses the dateline. \\
levels & path or list & \code{null} & HEALPix: interface depths, a level-table path or an inline list; needed in 3-D, and in 2-D it sets the node depths when the store has no bathymetry variable. \\
\bottomrule
\end{longtable}

\section*{\texttt{data} --- input and output}
\begin{longtable}{@{}>{\ttfamily\footnotesize}p{0.295\linewidth} >{\footnotesize}p{0.075\linewidth} >{\footnotesize}p{0.095\linewidth} >{\footnotesize}p{0.37\linewidth}@{}}
\toprule
\normalfont Key & Type & Default & Description \\
\midrule \endhead
ncinput & str & required & Velocity input: a file, a Zarr store, or (FESOM2) a directory of per-variable files; \code{fesom14}: a file, a directory or a glob, the files grouped by the variables they hold (\cref{sec:input-reader}). \\
ncinput\_v & str & \code{""} & Separate $v$-velocity file, if not in \conf{ncinput} (NEMO: overrides grid-V; HEALPix: refused). \\
ncinput\_w & str & \code{""} & Separate vertical-velocity file (3-D), if applicable (HEALPix: refused). \\
ncinput\_eta & str & \code{""} & HEALPix: separate sea-surface-height store. \\
inicloud & str & required & Release positions: rows of \code{lon lat [depth] [release\_sec]}, geographic (\cref{sec:life-release}). \\
particle\_offset & int & \code{null} & Global index of the file's first particle in a sharded release (\cref{sec:perf-shard}); \code{null} = the file's \code{\# kelp particle\_offset N} header, else 0. Keys every random draw. \\
output\_path & str & memory & Output; \code{.zarr} selects Zarr, otherwise NetCDF4 (\cref{ch:output}). \\
time\_averaging & str & guess & \code{mean} (shift timestamps to window midpoints) or \code{instant}; set explicitly (guess: \code{mean} for FESOM2 and HEALPix). \code{fesom14}: leave unset --- the reader resolves the record times itself and a set value is refused; the output records \code{"resolved\_by\_reader"} (\cref{sec:input-timeavg}). \\
time\_stamp & str & \code{null} & \code{fesom14} only (the three stamp keys are refused for every other type): where a stamp sits in its window, \code{start}, \code{centre} or \code{end}, or \code{instant} for snapshots. Given together with \conf{time\_window} or not at all; overrides the files' time bounds or output schedule with a warning. \\
time\_window & str & \code{null} & \code{fesom14} only: the averaging window (with \code{instant}, the snapshot cadence): \code{calendar\_month}, \code{calendar\_year}, \code{calendar\_day} or a duration \code{"<number> <s|min|h|d>"}, computed in the files' calendar. \\
time\_offset & str & \code{null} & \code{fesom14} only: a shift \code{"<$\pm$number> <s|min|h|d>"} added to the stamps, bounds and epoch of every file first, for raw output of a job that started after midnight; allowed alone. \\
\bottomrule
\end{longtable}

\section*{\texttt{time} --- window and realizations}
\begin{longtable}{@{}>{\ttfamily\footnotesize}p{0.295\linewidth} >{\footnotesize}p{0.075\linewidth} >{\footnotesize}p{0.095\linewidth} >{\footnotesize}p{0.37\linewidth}@{}}
\toprule
\normalfont Key & Type & Default & Description \\
\midrule \endhead
start\_time\_idx & int & \code{0} & First snapshot index to read. \\
n\_steps & int & --- & Number of integration steps (preferred way to set the length). \\
end\_time\_idx & int & --- & Legacy alternative to \conf{n\_steps}. \\
start\_date & str & --- & Start as an ISO date in the axis calendar; between snapshots only with \conf{temporal\_interp} (\cref{sec:input-window}). \\
duration\_days & float & --- & With \conf{start\_date}: run length in days, $\lfloor\text{days}\cdot86400/\dt\rfloor$ steps. \\
on\_exhausted & str & \code{stop} & At the end of the forcing: \code{stop} or \code{cycle} (wrap). \\
num\_runs & int & \code{1} & Independent realizations (ensemble; \cref{ch:diffusion}). \\
run\_interval & int & \code{0} & Snapshot-index offset between successive realizations. \\
n\_pulses & int & \code{1} & Re-release the seed set this many times (\cref{sec:life-release}). \\
pulse\_interval & float (s) & \code{0.0} & Interval between pulses (required when \conf{n\_pulses} $>1$). \\
\bottomrule
\end{longtable}

\section*{Command line}
The command \code{kelp run config.yaml} runs a configuration, \code{kelp info config.yaml} prints the mesh and
the time range without tracking, and \code{kelp --version} prints the engine version. The options of \code{kelp run} override the corresponding keys:

\begin{longtable}{@{}>{\ttfamily\footnotesize}p{0.295\linewidth} >{\footnotesize}p{0.58\linewidth}@{}}
\toprule
\normalfont Option & Effect \\
\midrule \endhead
-o, --output & Output path; overrides \conf{data.output\_path}. \\
--checkpoint-freq N & Overrides \conf{numerics.checkpoint\_freq}. \\
--temporal-interp & Switches the interpolated mode on; without the option the configuration decides. \\
--seed N & Overrides \conf{numerics.diffusion.seed}. \\
--flush-every N & Overrides \conf{numerics.output\_flush\_every}. \\
-v, --verbose & Per-step timing with the phase split of \cref{sec:perf-measure}. \\
\bottomrule
\end{longtable}
